\documentclass[11pt]{article}
\usepackage[T1]{fontenc}
\usepackage[a4paper,margin=1in]{geometry}
\usepackage{amsmath,amssymb,amsthm}
\usepackage{graphicx}
\usepackage[hidelinks]{hyperref}
\usepackage{microtype}

\newtheorem{theorem}{Theorem}[section]
\newtheorem{proposition}[theorem]{Proposition}
\newtheorem{remark}[theorem]{Remark}
\newcommand{\B}{\mathcal B}
\newcommand{\E}{\mathcal E}
\newcommand{\T}{\mathbb T}
\newcommand{\Tr}{\operatorname{Tr}}

\setlength{\parindent}{0pt}
\setlength{\parskip}{0.55em}

\title{Kernel Obstructions for Quantum Integration\\
and a Regular Scalar-Injective POVM}
\author{Run \texttt{fa\_banach\_001} --- lane 14}
\date{11 August 2026}

\begin{document}
\maketitle

\begin{center}
\fbox{\parbox{0.9\textwidth}{
\textbf{Status: candidate partial result, likely valid.}
Two broad density regimes force noninjectivity of the full quantum integration
map.  Conversely, an explicit nonatomic POVM with norm-continuous uniformly
coercive density has injective scalar integration, although its full quantum
integration map is noninjective and its range is nonconvex.  The general
injectivity question remains open only in a sharply isolated singular-density
regime.
}}
\end{center}

\section{Source problem and setup}

Plosker and Ramsey~\cite{source} define, for a POVM
$\nu:\Sigma\to\B(\mathcal H)$ and a full-rank density operator $\rho$,
the scalar measure
\[
 \mu=\nu_\rho,\qquad \mu(E)=\Tr(\rho\nu(E)),
\]
and the operator-valued Radon--Nikodym density
$D=d\nu/d\mu$, when it exists.  Their quantum integration map is
\[
 \E_\nu:L^\infty_{\mathcal H}(X,\nu)\longrightarrow\B(\mathcal H),
 \qquad \E_\nu(F)=\int_X F\,d\nu,
\]
characterized ultraweakly by
\begin{equation}\label{eq:integral}
 \Tr\!\left(s\E_\nu(F)\right)
 =\int_X\Tr\!\left(sD(x)^{1/2}F(x)D(x)^{1/2}\right)d\mu(x)
 \quad(s\in\mathcal T(\mathcal H)).
\end{equation}

On arXiv PDF page~7 they identify ``whether $\E_\nu$ is injective or not''
as a natural problem.  They then call scalar integration \emph{classically
non-injective} when, on every nonnull measurable subset, the restriction
$f\mapsto\E_\nu(fI)$ has a nonzero kernel.  This stronger local scalar
condition is a hypothesis of their operator-valued Lyapunov theorem.

\begin{figure}[ht]
\centering
\includegraphics[width=0.94\textwidth]{figures/open_problem_crop.png}
\caption{The injectivity problem and the classically non-injective condition
on source PDF page~7.}
\end{figure}

Throughout, $\mathcal H$ is separable, $\nu(X)=I$, $D$ exists, and $\nu$ is
nonatomic.  Since $\nu$ and $\mu=\nu_\rho$ are mutually absolutely
continuous, $\mu$ is nonatomic as well.

\section{Proof intuition}

Formula~\eqref{eq:integral} shows that quantum integration first sandwiches
$F(x)$ by $D(x)^{1/2}$ and then averages.  A pointwise kernel of $D$ therefore
gives an immediate kernel element.  At the opposite extreme, if $D$ is
uniformly invertible on any nonnull region, one can prescribe the sandwiched
integrand there: conjugate a fixed operator $A$ by $D^{-1/2}$ and multiply by
a nonzero scalar function of mean zero.

This leaves only one possible regime for injectivity: $D(x)$ must be
injective but arbitrarily small in some directions almost everywhere.  The
explicit circle example below shows why the scalar version does not settle
the full question.  Its off-diagonal blocks record every Fourier coefficient,
so scalar integration is injective; nevertheless uniform invertibility lets
the operator-valued cancellation above work.

\section{Two kernel criteria}

\begin{theorem}[Pointwise-kernel criterion]\label{thm:kernel}
Suppose
\[
 K=\{x\in X:\ker D(x)\ne\{0\}\}
\]
has positive $\mu$-measure.  Then $\E_\nu$ is not injective.
\end{theorem}

\begin{proof}
Let $P(x)=\mathbf 1_{\{0\}}(D(x))$ on $K$, and put $P(x)=0$ off $K$.
This is a weakly measurable projection field: for example it is the strong
pointwise limit of $\mathbf 1_K(x)(I+nD(x))^{-1}$.  It is essentially bounded,
and it is nonzero in $L^\infty_{\mathcal H}(X,\nu)$ because $\mu(K)>0$ and
$\nu\sim\mu$.  Pointwise,
\[
 D(x)^{1/2}P(x)D(x)^{1/2}=0.
\]
Equation~\eqref{eq:integral} therefore gives $\E_\nu(P)=0$.
\end{proof}

\begin{theorem}[Local-coercivity criterion]\label{thm:coercive}
Suppose there are $E\in\Sigma$ and $c>0$ such that $\mu(E)>0$ and
\[
 D(x)\ge cI\qquad\text{for $\mu$-almost every }x\in E.
\]
Then $\E_\nu$ is not injective.  More precisely, for every nonzero
$A\in\B(\mathcal H)$ there is a nonzero $F\in
L^\infty_{\mathcal H}(X,\nu)$ with $\E_\nu(F)=0$.
\end{theorem}

\begin{proof}
By nonatomicity, split $E$ into measurable sets $E_1,E_2$ of positive
$\mu$-measure and set
\[
 g=\mu(E_2)\mathbf 1_{E_1}-\mu(E_1)\mathbf 1_{E_2}.
\]
Then $g$ is nonzero, supported in $E$, and $\int_Eg\,d\mu=0$.  Functional
calculus gives $\|D(x)^{-1/2}\|\le c^{-1/2}$ on $E$.  Define
\[
 F(x)=
 \begin{cases}
 g(x)D(x)^{-1/2}AD(x)^{-1/2},&x\in E,\\
 0,&x\notin E.
 \end{cases}
\]
This is measurable and essentially bounded, is nonzero, and satisfies
\[
 D(x)^{1/2}F(x)D(x)^{1/2}=g(x)A.
\]
Thus $\E_\nu(F)=A\int_Eg\,d\mu=0$.
\end{proof}

\begin{proposition}[Necessary shape of an injective density]
If $\E_\nu$ is injective, then for $\mu$-almost every $x$, the positive
operator $D(x)$ is injective and has dense, nonclosed range.
\end{proposition}

\begin{proof}
Theorem~\ref{thm:kernel} gives $\ker D(x)=0$ almost everywhere, hence
$\overline{\operatorname{ran}D(x)}=\mathcal H$.  If $D(x)$ were bounded
below on a positive-measure set, then one of the measurable sets
$\{x:D(x)\ge n^{-1}I\}$ would have positive measure, contradicting
Theorem~\ref{thm:coercive}.  Thus $D(x)$ is not bounded below almost
everywhere.  An injective positive operator with closed range is bounded
below, so its range here is nonclosed.
\end{proof}

\section{A regular scalar-injective example}

\begin{theorem}[Fourier POVM]\label{thm:fourier}
There is a nonatomic quantum probability measure $\nu$ on the circle with a
norm-continuous Radon--Nikodym density satisfying
\[
 \tfrac34 I\le D(z)\le\tfrac54 I,
\]
such that:
\begin{enumerate}
\item scalar integration $T(f)=\E_\nu(fI)$ is injective on
      $L^\infty(\T)$, so $\nu$ is not classically non-injective;
\item the full quantum integration map $\E_\nu$ is noninjective;
\item the range $\{\nu(E):E\in\Sigma\}$ is not convex.
\end{enumerate}
\end{theorem}

\begin{proof}
Let
$\mathcal H=\bigoplus_{k\ge1}\mathbb C^2$, with basis
$e_{k,0},e_{k,1}$ in the $k$th block.  Put $a_k=2^{-k-1}$ and, for
$z\in\T$, define the block diagonal operator
\[
 D(z)|_{\mathbb C^2_k}=
 \begin{pmatrix}
 1&a_k z^k\\
 a_k\overline z^{,k}&1
 \end{pmatrix}.
\]
Each block has eigenvalues $1\pm a_k$, so
$\frac34I\le D(z)\le\frac54I$.  Since $a_k\to0$, the block tail is
uniformly small and the finitely many leading blocks are continuous; hence
$z\mapsto D(z)$ is norm-continuous.  With $m$ normalized Haar measure,
\[
 \int_\T D(z)\,dm(z)=I.
\]
Define $\nu(E)=\int_E D(z)\,dm(z)$.

Take the full-rank density operator
\[
 \rho|_{\mathbb C^2_k}=2^{-k-1}I_2.
\]
Its trace is $\sum_{k\ge1}2\cdot2^{-k-1}=1$, and the diagonal form of
$\rho$ gives $\Tr(\rho D(z))=1$.  Therefore $\nu_\rho=m$ and
$D=d\nu/d\nu_\rho$.  Mutual absolute continuity with nonatomic $m$ shows
that $\nu$ is nonatomic.

For $f\in L^\infty(\T)$, the diagonal entries of $T(f)$ equal
$\int f\,dm$, while the two off-diagonal entries in block $k$ equal
\[
 a_k\int_\T f(z)z^k\,dm(z),
 \qquad
 a_k\int_\T f(z)\overline z^{,k}\,dm(z).
\]
Thus $T(f)=0$ forces every Fourier coefficient of $f$ to vanish.  Since
$f\in L^1(\T)$, Fourier uniqueness gives $f=0$ almost everywhere.  Hence
$T$ is injective, and so is its restriction to functions supported on any
measurable subset.

On the other hand, Theorem~\ref{thm:coercive} applies globally, proving that
$\E_\nu$ is noninjective.  Explicitly, for any nonzero $A\in\B(\mathcal H)$,
\[
 F(z)=zD(z)^{-1/2}AD(z)^{-1/2}
\]
is bounded and nonzero, while $\E_\nu(F)=A\int_\T z\,dm(z)=0$.

Finally, if the range of $\nu$ were convex, it would contain the midpoint
$I/2$ of $0=\nu(\varnothing)$ and $I=\nu(\T)$.  Thus some measurable $E$
would satisfy
\[
 T(\mathbf1_E)=\nu(E)=\tfrac12I=T(\tfrac12).
\]
Injectivity of $T$ would force $\mathbf1_E=1/2$ almost everywhere, which is
impossible.  Hence the range is not convex.
\end{proof}

\section{Limitations, upgrade attempts, and novelty}

The results do not decide whether an injective $\E_\nu$ exists in the only
remaining regime: $D(x)$ injective with dense nonclosed range almost
everywhere.  Six focused upgrade attempts included inverse-density kernels,
measurable kernel projections, a commutator cancellation formula, a
Stinespring compression reformulation, a predual density reformulation, and
a search for an injective Fourier model.  The commutator and inversion routes
become unbounded exactly in the remaining regime; the functional-analytic
reformulations supplied no contradiction.  The Fourier search instead
produced Theorem~\ref{thm:fourier}, which sharply separates scalar from full
injectivity.

A bounded search through 11 August 2026 checked the run indexes, local arXiv
corpus, exact source title and id, and web/arXiv queries combining
``classically non-injective'', ``POVM integration injective'',
``operator-valued Lyapunov'', and ``Radon--Nikodym density''.  It found the
source and general informationally complete POVM literature, but no later
resolution of the stated injectivity problem and no occurrence of the block
Fourier construction above.  The novelty verdict is therefore provisional
pending expert review.

\section{Human-review recommendation}

The partial result is recommended for expert review.  The main checks are:
measurability of the kernel projection in Theorem~\ref{thm:kernel}, domain
membership of the inverse-density kernel in Theorem~\ref{thm:coercive}, and
the Fourier uniqueness identification in Theorem~\ref{thm:fourier}.  The
packet should not be promoted to a full solution unless the singular
injective-density regime is settled.

\begin{thebibliography}{9}
\bibitem{source}
S. Plosker and C. Ramsey,
\emph{An operator-valued Lyapunov theorem},
J. Math. Anal. Appl. \textbf{469} (2019), 117--125;
arXiv:1801.03109.
\end{thebibliography}

\end{document}
